% ============================================================
% ch12.tex —— 第 12 章（选读）不可能的作图：代数给几何判死刑
% ============================================================
\chapter{不可能的作图：代数给几何判死刑}\label{ch:zuotu}

\noindent\textbf{本章为选读}（难度 $\bigstar\bigstar\bigstar$），跳过它不影响后续任何一章。但如果你留下，你将见证数学史上最气势恢弘的一次跨界执法：\textbf{用代数，给几何问题判死刑}——证明某件事\emph{永远不可能}。"不可能"的证明是数学的最高形式之一：一亿次失败不算数，必须从原理上封死。

\section{两千年的悬案}

规矩先立好。只许用两件工具：\textbf{无刻度直尺}（只能连线，不能量长度）和\textbf{圆规}（只能画圆或截取线段，不能当量角器）。从一条已知线段（长度记为 $1$）出发作图。古希腊人提出了三大难题，两千三百年无人破解：

\begin{itemize}
\item \textbf{化圆为方}：作一个正方形，面积等于已知圆；
\item \textbf{倍立方}：作一个立方体，体积是已知立方体的两倍（传说雅典瘟疫时，神谕要求把祭坛加倍而不许改变形状）；
\item \textbf{三等分任意角}：把任意角精确分成三等份。
\end{itemize}

无数次"成功"的报告后来全被发现是作弊（偷用了刻度、或只是近似）。僵局的根源在于所有人的思路都是"\textbf{想办法作出来}"。破局只需要一句反问：\textbf{要是它们根本作不出来呢？}要证明"作不出来"，硬试无用——必须从根上说明\textbf{尺规这套工具，天产能造出哪些长度、造不出哪些}。这就要给"尺规的能力"画一条生产线流程图——而画图的工具，是代数。

\section{尺规到底会干什么}

拆开看，尺规的全部本领只有两个动作：
\begin{enumerate}
\item 过两个已知的点画直线；以已知点为圆心、已知两点距离为半径画圆；
\item 取\textbf{交点}：直线与直线、直线与圆、圆与圆的交点，成为新的已知点。
\end{enumerate}
没有别的了。一切的长度都从交点的坐标里长出来。现在把两个动作翻译成代数（初中坐标系的存货刚好够用）：

\paragraph{直线与直线相交。}两条直线 $a_1 x + b_1 y = c_1$、$a_2 x + b_2 y = c_2$（系数由已知点写出），联立解出交点——加减消元，只用到\textbf{加减乘除}。

\paragraph{直线与圆相交。}圆 $(x-u)^2 + (y-v)^2 = r^2$，代入直线 $y = mx + k$，得到关于 $x$ 的\textbf{一元二次方程}。求根公式：$x = \dfrac{-b \pm \sqrt{b^2 - 4ac}}{2a}$——加减乘除，外加\textbf{一次开平方}。

\paragraph{圆与圆相交。}两圆方程相减（二次项 $x^2, y^2$ 恰好抵消），得到一条直线（两圆的"公共弦"）——化归为上一情形，还是\textbf{加减乘除加开平方}。

结论打包：

\begin{dingli}[可作图数的生产线]
从长度 $1$ 出发，尺规能作出的每一个长度，都可以由 $1$ 经过\textbf{有限次加减乘除与开平方}得到。它们住在这样一座塔里：
\[
\mathbb{Q} \;\to\; \mathbb{Q}(\sqrt2) \;\to\; \mathbb{Q}\bigl(\sqrt2,\ \sqrt{2+\sqrt3}\bigr) \;\to\; \cdots
\]
塔的第一层是全体分数（加减乘除造得出）；每升一层，只许\textbf{多开一次平方根}。\textbf{尺规是台只能开平方的机器}——开立方、开五次方，统统超纲。
\end{dingli}

（每步画图只添一次平方根，但画很多步可以一层层累积——所以是"塔"而不是"一步到位"。塔有无限多层，但每个具体作图只能爬有限层。）

\section{判死刑：三个案子逐一过堂}

\paragraph{案子一：三等分 $60^\circ$ 角。}它等价于作出长度 $\cos 20^\circ$（作出这条直角边就能拼出 $20^\circ$ 角）。借一条三角恒等式（三倍角公式，高中内容，此处只用结论）：
\[
\cos 3\theta = 4\cos^3\theta - 3\cos\theta .
\]
令 $\theta = 20^\circ$：左边 $\cos 60^\circ = \tfrac12$。记 $c = \cos 20^\circ$，得
\[
4c^3 - 3c = \tfrac12 \quad\Longleftrightarrow\quad 8c^3 - 6c - 1 = 0 . \tag{$*$}
\]
审问 $(*)$：它有没有分数根？三次方程若有分数根 $\tfrac pq$（既约），$p$ 必整除常数项 $-1$、$q$ 必整除首项 $8$——候选只有 $\pm1, \pm\tfrac12, \pm\tfrac18, \pm\tfrac14$。逐个代入（动手！）：
\[
\begin{array}{ll}
c = 1: & 8 - 6 - 1 = 1 \ne 0 \\
c = -1: & -8 + 6 - 1 = -3 \ne 0 \\
c = \tfrac12: & 1 - 3 - 1 = -3 \ne 0 \\
c = -\tfrac12: & -1 + 3 - 1 = 1 \ne 0 \\
\end{array}
\]
（$\pm\tfrac14, \pm\tfrac18$ 代入同样不为零。）全灭——\textbf{$(*)$ 没有分数根}。而"从分数出发经过加减乘除得到的一切数"恰好是分数（分数对四则运算封闭，第 10 章），所以 $c$ \textbf{不在塔的第一层}。它能进塔吗？塔的每一层相对下一层只开平方根——可以证明（本节末尾的技术框）这使"塔内居民满足的整系数方程次数"只能是 $1, 2, 4, 8, \dots$（$2$ 的幂）。而 $c$ 满足的 $(*)$ 是三次且没有更低次的分数关系——次数 $3$，\textbf{不是 $2$ 的幂}。\textbf{死刑：三等分 $60^\circ$，尺规永不可能。}

\paragraph{案子二：倍立方。}目标棱长 $x$ 满足 $x^3 = 2$。同款审问：$x^3 - 2 = 0$ 没有分数根（候选 $\pm1, \pm2$ 全灭），次数 $3$ 不是 $2$ 的幂——\textbf{死刑}。（顺便，$\sqrt[3]{2}$ 用标了刻度的尺子当然量得出来——判的是"尺规净土内的不可能"，不是"数不存在"。）

\paragraph{案子三：化圆为方。}目标正方形边长 $\sqrt\pi$。这一次连方程都列不出：$\pi$ 根本不是任何整系数多项式的根（1882 年林德曼证明 $\pi$ 是\textbf{超越数}——不满足任何代数方程）。塔的户口本（代数数）上查无此人，\textbf{死刑}，且是最彻底的一种：连"次数"都谈不上。

三大难题，两千三百年，三张死刑判决书，依据只有一句话：\textbf{你要的数，生产的机器造不出}。数学史上第一次，"不可能"本身成了定理。

\begin{gongju}
\textbf{技术框：为什么塔内居民的次数只能是 $2$ 的幂？}粗算一笔：第一层是分数（次数 $1$）。在已有的数库里添一个平方根 $\sqrt{m}$，新库里的每个数形如 $a + b\sqrt m$（$a, b$ 是旧库的数）——新库作为"旧库上的向量空间"是 $2$ 维的（基底 $1, \sqrt m$）。塔每升一层，维度至多翻倍：$1 \to 2 \to 4 \to 8 \to \cdots$。而"一个数满足的最低次整系数方程的次数"$\le$ 它所在层的维度（线性代数的小推论）。所以塔内居民的次数至多是 $1, 2, 4, 8, \dots$——$3$ 永远挤不进去。这套"数库逐层加盖"的思想，正式名字叫\textbf{域扩张}，它是伽罗瓦理论的地基。
\end{gongju}

\begin{cidian}
土名"塔内居民"$\to$\textbf{可作图数}；塔本身 $\to$ \textbf{二次扩张塔}；"满足整系数方程的数"$\to$\textbf{代数数}；"不满足任何整系数方程的数"$\to$\textbf{超越数}（$\pi$、$\mathrm{e}$ 都是，$\mathrm{e}$ 的超越性 1873 年由埃尔米特证明）。
\end{cidian}

\begin{lishi}
\textbf{埃瓦里斯特·伽罗瓦}（1811--1832），把本章思想推到极致的人。两次报考巴黎综合理工学校落榜（传说一次是把黑板擦扔到考官头上——真假存疑，但他的暴烈属实）；三篇论文两次被柯西、泊松等院士遗失或退回；因参加共和派活动入狱；1832 年 5 月 29 日，决斗前夜，他通宵给好友舍瓦利耶写信，把数学遗嘱（后世的"伽罗瓦理论"）匆匆写就，页边反复出现"我没有时间了"。次日中弹，两日后去世，差半周满二十一岁。他留下的思想是：\textbf{给每个方程配一个群（根的对称群），方程可不可解（有没有根号解），取决于这个群能不能切成逐层"素数阶"的塔}——可解群的概念由此诞生。拉格朗日的盘子（第 9 章）在他手中成为解剖刀；五次方程没有求根公式（阿贝尔 1824 年先证）只是这套理论的一个特例。十四年后手稿被刘维尔整理发表，世界才意识到那封信里装着什么。
\end{lishi}

\begin{dongshou}
\begin{itemize}
  \yisi 用塔的规则解释：正三角形、正方形、正六边形可以尺规作出（余弦 $\cos 60^\circ = \tfrac12$、$\cos 90^\circ = 0$、$\cos 30^\circ = \tfrac{\sqrt3}{2}$——都只开一层平方根，塔内居民）；正九边形不行（要 $\cos 40^\circ$，而 $40^\circ$ 三等分了 $120^\circ$……准确说 $\cos 40^\circ$ 满足三次方程 $8c^3 - 6c = \cos 120^\circ = -\tfrac12$，与案子一同款）。
  \yier 把案子一的验根做全：$8c^3 - 6c - 1$ 在 $c = \pm\tfrac14, \pm\tfrac18$ 时的值，逐个算出（$c = \tfrac14$：$8\cdot\tfrac1{64} - \tfrac32 - 1 = \tfrac18 - \tfrac52 = -\tfrac{19}{8} \ne 0$，等等），确认无分数根。
  \yier 倍立方案的完整验根：$x^3 - 2$ 在 $\pm1, \pm2$ 处的值。（$1-2 = -1$，$-1-2 = -3$，$8-2 = 6$，$-8-2 = -10$，全非零 \checkmark。）
\end{itemize}
\end{dongshou}

\begin{shaonao}
\textbf{高斯的十七边形}。正 $n$ 边形可尺规作图 $\iff$ $n$ 是"二的一个幂乘若干个不同的费马素数"（$n = 2^k p_1\cdots p_s$，$p_i$ 形如 $2^{2^m}+1$ 的素数）。$17 = 2^4 + 1$ 是费马素数，所以\textbf{正十七边形可以尺规作出}——十九岁的高斯在 1796 年 3 月 30 日早晨证明了这一点，兴奋到从此立志数学（此前他在语言与数学之间摇摆）；这个发现也铭刻在了他的墓碑上。你可以亲手验证链条的入口：$\cos\dfrac{360^\circ}{17}$ 满足的方程次数是 $16 = 2^4$——$2$ 的幂，塔内合法居民（它的四个"中间方程"各为二次，一层层开上去）。对照正七边形：$\cos\dfrac{360^\circ}{7}$ 的次数是 $3$——死刑。\textbf{能作与不能作的分界线，就是一条次数的分数线}——你刚刚用初等工具走完了它的全部逻辑。
\end{shaonao}

\begin{chengshi}
本章三处从简：三倍角公式直接引用（高中内容）；"三次方程无分数根 $\Rightarrow$ 次数恰为 $3$"（完整说法是"$8x^3-6x-1$ 在有理数上不可约"，我们用验根代替了不可约性论证——对三次方程，"没有有理根"确实等价于"不可约"，这一步其实是干净的）；"塔内次数是 $2$ 的幂"的技术框是导引版（维度论证严格化在大学线性代数之后）。另外注意"可作图 $\Rightarrow$ 次数是 $2$ 的幂"只是必要条件——次数是 $2$ 的幂未必可作图（高斯判据才是完整的充要条件）。三张死刑判决书本身（1837 年 Wantzel 两案、1882 年 Lindemann 一案）则是数学史上无争议的定案。
\end{chengshi}

篇二完。装备清单：群、群表、同构、子群、陪集、拉格朗日定理、环、域、零因子、分数流水线、多项式环、余数定理、可作图数。下一篇，我们去驯服数学里最狂野的概念——\textbf{无穷}。
